% Work Package 1: Conformal Training with an Encoder on Frozen Step Embeddings
\documentclass[10pt,a4paper]{article}
% Shared clean typography and visual layout for WP1 & WP2
\usepackage[margin=1.85cm,top=1.5cm,bottom=1.55cm,footskip=0.75cm]{geometry}
\usepackage[T1]{fontenc}
\usepackage[utf8]{inputenc}
\usepackage{lmodern}
\usepackage{microtype}
\usepackage{amsmath,amssymb,amsthm}
\usepackage{enumitem}
\usepackage{booktabs,tabularx,array}
\usepackage{xcolor}
\usepackage{titlesec}
\usepackage[skins]{tcolorbox}
\usepackage[font=small,labelfont={bf,sf},skip=4pt]{caption}
\usepackage[hidelinks]{hyperref}

% ---------- calm palette ----------
\definecolor{ink}{HTML}{1F2937}
\definecolor{muted}{HTML}{52606D}
\definecolor{navy}{HTML}{1E3A5F}
\definecolor{teal}{HTML}{1F6E7A}
\definecolor{tint}{HTML}{F6F8FA}
\definecolor{tint2}{HTML}{EEF4F6}
\definecolor{rulegray}{HTML}{CBD5E1}

\color{ink}
\linespread{1.13}
\setlength{\parskip}{3.5pt}
\setlength{\parindent}{0pt}
\setlength{\textfloatsep}{8pt plus 2pt}
\setlength{\floatsep}{6pt plus 2pt}
\setlength{\intextsep}{7pt plus 2pt}
\setlength{\abovedisplayskip}{5pt plus 1pt}
\setlength{\belowdisplayskip}{5pt plus 1pt}
\setlength{\emergencystretch}{2em}
\setlist{nosep,leftmargin=1.35em,itemsep=2pt}
\pagestyle{plain}

% ---------- section headings ----------
\titleformat{\section}
  {\sffamily\bfseries\large\color{navy}}
  {\thesection}{0.55em}{}
  [\vspace{-6pt}\textcolor{rulegray}{\rule{\linewidth}{0.5pt}}]
\titlespacing*{\section}{0pt}{8pt plus 1.5pt minus 1pt}{3pt}

\titleformat{\subsection}
  {\sffamily\bfseries\normalsize\color{ink}}
  {\thesubsection}{0.5em}{}
\titlespacing*{\subsection}{0pt}{5.5pt plus 1pt minus 1pt}{2pt}

\newcolumntype{L}{>{\raggedright\arraybackslash}X}

% ---------- reusable visual blocks ----------
\newcommand{\wpheader}[4]{%
  \begin{tcolorbox}[
    enhanced, colback=navy, colframe=navy, arc=3pt, boxrule=0pt,
    left=10pt, right=10pt, top=8pt, bottom=8pt
  ]
  {\sffamily\footnotesize\color{white!75}\MakeUppercase{#1}}\\[2pt]
  {\sffamily\bfseries\Large\color{white}#2}\\[3pt]
  {\sffamily\small\color{white!88}#3}
  \end{tcolorbox}
  \vspace{-3pt}
  \begin{tcolorbox}[
    enhanced, colback=tint, colframe=rulegray, arc=2pt, boxrule=0.5pt,
    left=8pt, right=8pt, top=5pt, bottom=5pt
  ]
  \small\sffamily #4
  \end{tcolorbox}
  \vspace{2pt}
}

\tcbset{
  algbox/.style={
    enhanced,
    colback=tint,
    colframe=rulegray,
    borderline west={2.5pt}{0pt}{navy},
    boxrule=0.5pt,
    arc=2pt,
    left=8pt,right=7pt,top=6pt,bottom=6pt,
    fonttitle=\sffamily\bfseries\small\color{navy},
    colbacktitle=tint,
    coltitle=navy,
    boxed title style={empty},
    title style={left=0pt}
  },
  critbox/.style={
    enhanced,
    colback=tint2,
    colframe=teal!35,
    borderline west={2.5pt}{0pt}{teal},
    boxrule=0.5pt,
    arc=2pt,
    left=8pt,right=7pt,top=5pt,bottom=5pt,
    fonttitle=\sffamily\bfseries\small\color{teal},
    colbacktitle=tint2
  }
}

% ---------- shared math macros ----------
\newcommand{\Prob}{\mathbb{P}}
\newcommand{\E}{\mathbb{E}}
\newcommand{\R}{\mathbb{R}}
\newcommand{\ind}{\mathbf{1}}
\newcommand{\C}{C}
\newcommand{\Qhat}{\widehat Q}
\newcommand{\Smin}{S^{\min}}
\newcommand{\Dtr}{D_{\mathrm{train}}}
\newcommand{\Dcal}{D_{\mathrm{cal}}}
\newcommand{\Dtest}{D_{\mathrm{test}}}
\newcommand{\Dval}{D_{\mathrm{val}}}
\newcommand{\tagpill}[1]{\textcolor{navy}{\sffamily\bfseries #1}}


\newcommand{\Lcls}{\widehat{\mathcal L}_{\mathrm{cls}}}
\newcommand{\Lsize}{\widehat{\mathcal L}_{\mathrm{size}}}

\begin{document}

\wpheader
  {Work Package 1 \textperiodcentered\ Implementation Spec}
  {Conformal Training with an Encoder on Frozen Step Embeddings}
  {Supervised DPSM conformal training over frozen LLM step representations, evaluated with vanilla split CP}
  {\tagpill{Backbone:} Qwen3-8B (frozen) \ \textcolor{rulegray}{|}\ 
   \tagpill{Scorer:} 2-layer Transformer \ \textcolor{rulegray}{|}\ 
   \tagpill{Score:} $S=1-\pi_t$ \ \textcolor{rulegray}{|}\ 
   \tagpill{Train:} DPSM ($\alpha=0.1$) \ \textcolor{rulegray}{|}\ 
   \tagpill{Eval:} $\alpha\in\{0.05,0.1,0.2\}$}

% =====================================================================
\section{Implementation Steps}

\subsection{Step 1: Cache Frozen Step Embeddings}
\begin{enumerate}
\item Load frozen \textbf{Qwen3-8B} (\texttt{bfloat16}, \texttt{eval()}, no grad). Format each trace $x=(s_1,\dots,s_T)$ as task + ordered steps and record token spans $I_t$.
\item Run one forward pass per trace (for traces exceeding context length, use $50\%$-overlapping windows and take each step from the window with the longest left context).
\item Extract hidden states at candidate layers $\ell\in\{25\%,50\%,75\%,100\%\}$ depth (pick $\ell$ on $\Dval$) and mean-pool per step:
\[
h_t = \frac{1}{|I_t|}\sum_{j\in I_t} H_j^{(\ell)} \in \R^{d_{\mathrm{llm}}}, \qquad t\in[T].
\]
\item Save $\{h_1,\dots,h_T\}$, $T$, and $Y$ to disk for all splits and benchmarks.
\end{enumerate}

\subsection{Step 2: Step-Scorer Architecture}
Define $f_\theta$ with trainable parameters $\theta=(W, \phi, w)$ and hidden width $d_{\mathrm{enc}}=256$:
\begin{equation}
u_t = \mathrm{LN}(W h_t), \qquad
(v_1,\dots,v_T) = \mathrm{Enc}_\phi(u_1,\dots,u_T), \qquad
\pi_t(x) = \frac{\exp(w^\top v_t)}{\sum_{t'=1}^{T}\exp(w^\top v_{t'})}. \label{eq:wp1-arch}
\end{equation}
Here $\mathrm{Enc}_\phi$ is a 2-layer bidirectional Transformer encoder (4 heads, FFN $1024$, dropout $0.1$, key-padding mask over valid steps $t\in[T]$) applying the backbone's positional encoding over step positions $1,\dots,T$ (RoPE for Qwen3).

\subsection{Step 3: Conformal Training}
Split $\Dtr\setminus\Dval$ once by task into disjoint halves $D_1, D_2$ and run Algorithm~1 (see Appendix~\ref{app:conftr} for a ConfTr alternative).

\begin{tcolorbox}[algbox,title={Algorithm 1 \textperiodcentered\ DPSM Conformal Training for Step-Embedding Encoder}]
\small
\textbf{Inputs:} $D_1, D_2$; target $\alpha=0.1$; weights/rates $\lambda, \tau, \eta, \gamma$; batch size $b$; steps $K$; $\varepsilon=10^{-6}$. \quad \textbf{Init:} random $\theta_0$, $q_0=0.5$.
\begin{enumerate}[leftmargin=1.4em,itemsep=2pt]
\item \textbf{For} $k=1,\dots,K$:
  \begin{enumerate}[label=\arabic{enumi}.\arabic*.,leftmargin=1.6em,itemsep=1.5pt]
  \item Sample batches $B^1\subset D_1$ and $B^2\subset D_2$ of size $b$.
  \item \textbf{Classification loss ($B^1$):} $\displaystyle \Lcls(\theta) = \frac{1}{|B^1|}\sum_{i\in B^1} -\log\!\Big(\sum_{t\in Y_i}\pi_t(x_i)+\varepsilon\Big)$.
  \item \textbf{Relative soft set-size loss ($B^2$, \texttt{stop\_grad($q$)}):} $\displaystyle \Lsize(\theta,q) = \frac{1}{|B^2|}\sum_{j\in B^2}\frac{1}{T_j}\sum_{t=1}^{T_j}\sigma\!\left(\frac{q - 1 + \pi_t(x_j)}{\tau}\right)$.
  \item \textbf{Encoder step:} $\theta \leftarrow \mathrm{AdamW}\big(\theta,\; \nabla_\theta[\Lcls(\theta)+\lambda\Lsize(\theta,q)],\; \eta\big)$.
  \item \textbf{Threshold step ($B^1$, no grad to $\theta$):} with pre-update $\Smin_i = 1 - \max_{t\in Y_i}\pi_t(x_i)$,
  \[
  q \leftarrow \mathrm{clip}_{[0,1]}\!\left(q - \gamma\left(\alpha - \frac{\#\{i\in B^1 : \Smin_i \ge q\}}{|B^1|}\right)\right).
  \]
  \end{enumerate}
\item \textbf{Output:} Discard $q$; select best checkpoint $\theta^\star$ on $\Dval$ and calibrate $\Qhat(\alpha)$ on $\Dcal$ (Section~\ref{sec:vanilla-cp}).
\end{enumerate}
\end{tcolorbox}

\textbf{Sweep grid (select on $\Dval$):} AdamW ($b=32$, 50 epochs, weight decay $10^{-2}$, cosine schedule); $\eta\in\{10^{-4},3\!\cdot\!10^{-4},10^{-3}\}$; $\lambda\in\{0.1,0.5,1,2,5\}$; $\tau\in\{0.01,0.05,0.1\}$; $\gamma\in\{10^{-3},5\!\cdot\!10^{-3},10^{-2},5\!\cdot\!10^{-2}\}$; 5 seeds at best config.

% =====================================================================
\newpage
% Shared Vanilla Split CP, Datasets, and Evaluation Protocol (WP1 & WP2)

\section{Shared Vanilla Split CP and Evaluation Setup} \label{sec:shared}

\subsection{Vanilla Split Conformal Prediction} \label{sec:vanilla-cp}

Given any trained scorer outputting $\pi_\theta(x)=(\pi_1,\dots,\pi_T)\in\Delta^{T-1}$ and a held-out calibration set $\Dcal=\{(x_i,Y_i)\}_{i=1}^m$ (each $Y_i\subseteq[T]$ non-empty), run three steps at target miscoverage $\alpha\in(0,1)$:
\begin{enumerate}
\item \textbf{Score:} $S_\theta(x,t) = 1 - \pi_t(x)$ for each step $t\in[T]$; trajectory calibration score $\Smin_i = \min_{t\in Y_i} S_\theta(x_i,t) = 1 - \max_{t\in Y_i}\pi_t(x_i)$.
\item \textbf{Calibrate:} Discard any training threshold $q$. Set $k=\lceil(1-\alpha)(m+1)\rceil$ and take $\Qhat(\alpha)$ as the $k$-th smallest value of $\{\Smin_1,\dots,\Smin_m\}$ (or $+\infty$ if $k=m+1$).
\item \textbf{Predict:} For test trajectory $x$, return $\C_\theta(x) = \{\,t\in[T] : S_\theta(x,t) \le \Qhat(\alpha)\,\}$.
\end{enumerate}

\subsection{Datasets, Splits, and Metrics} \label{sec:datasets}

\smallskip
{\small\renewcommand{\arraystretch}{1.06}%
\noindent\begin{tabularx}{\linewidth}{@{}>{\raggedright\arraybackslash}p{2.8cm} L@{}}\toprule
\sffamily\bfseries Element & \sffamily\bfseries Shared Specification \\\midrule
\textbf{Backbone} & Open-weight \textbf{Qwen3-8B} for both WP1 (frozen hidden states) and WP2 (LoRA fine-tuning). \\
\textbf{Synthetic corpus} & Failed Magentic-One \& CaptainAgent trajectories with counterfactual error sets $Y\subseteq[T]$ ($Y\neq\emptyset$). Exclude all GAIA val and AssistantBench dev tasks to prevent benchmark leakage. \\
\textbf{Splits} & Task-disjoint $\Dtr, \Dcal, \Dtest$. Reserve task-disjoint $\Dval\subset\Dtr$ for selection; halve $\Dtr\setminus\Dval$ once by task into $D_1, D_2$. Use identical label verification on $\Dcal,\Dtest$ and $\le 1$ trajectory per task in $\Dcal$. \\
\textbf{External benchmarks} & \textbf{Who\&When} (184, $|Y|=1$), \textbf{TraceElephant} (176, $|Y|=1$), \textbf{AgentRx} (44, multi-step), \textbf{MP-Bench} (289; union \& 2-of-3). Report side by side (never pooled) under: (i) \emph{Recalibrated} (200$\times$ random 50/50 cal/test splits; $\alpha\in\{0.1,0.2\}$, plus $0.05$ where $m\ge 40$), (ii) \emph{Zero-shot transfer} (calibrate on synthetic $\Dcal$; report coverage gap as shift diagnostic). \\
\textbf{Model selection} & Pick checkpoint/hyperparameters on $\Dval$ (internal cal/eval split) by lowest calibrated $\widehat{\mathrm{Ineff}}_{\mathrm{rel}}$ at $\alpha=0.1$, identically for every arm. Never touch $\Dcal$ or $\Dtest$. \\
\textbf{Headline metrics} & Evaluate at $\alpha\in\{0.05,0.1,0.2\}$ (WP1: 5 seeds; WP2: 3 seeds):
$\widehat{\mathrm{Cov}} = \frac{1}{N}\sum_j \ind[\C(x_j)\cap Y_j\neq\emptyset]$,
$\widehat{\mathrm{Ineff}} = \frac{1}{N}\sum_j |\C(x_j)|$,
$\widehat{\mathrm{Ineff}}_{\mathrm{rel}} = \frac{1}{N}\sum_j |\C(x_j)|/T_j$.
Check $\widehat{\mathrm{Cov}}$ against $\mathrm{Beta}(k,m+1-k)$ + test bootstrap CI. \\
\textbf{Guardrails} & Report empty-set rate, full-set rate, coverage by $T$-tercile and agent system, and top-1 hit rate. \\\bottomrule
\end{tabularx}}
\smallskip


% =====================================================================
\section{Required Baselines and Claim Criteria} \label{sec:baselines}

\smallskip
{\small\renewcommand{\arraystretch}{1.04}%
\noindent\begin{tabularx}{\linewidth}{@{}>{\raggedright\arraybackslash}p{3.1cm} L L@{}}\toprule
\sffamily\bfseries Baseline & \sffamily\bfseries Specification & \sffamily\bfseries Claim Isolated \\\midrule
\textbf{B1. CE (full $\Dtr$)} & Encoder trained with $\Lcls$ on all of $\Dtr\setminus\Dval$ + vanilla split CP. & Primary: does DPSM beat standard full-data CE? \\
\textbf{B2. CE ($D_1$ only)} & Same as B1, trained on $D_1$ only. & Equal-label control for halving $\Lcls$ to $D_1$. \\
\textbf{B3. CE + fixed $q_{\mathrm{fix}}$} & Train with $\Lcls+\lambda\Lsize(\theta,q_{\mathrm{fix}})$; constant $q_{\mathrm{fix}}$ tuned on $\Dval$. & Value of learning $(1-\alpha)$-quantile $q$ vs.\ fixed threshold. \\
\textbf{B4. Temp.\ scaling} & Tune $T_{\mathrm{temp}}>0$ on B1 logits ($\mathrm{softmax}(\ell/T_{\mathrm{temp}})$) on $\Dval$ + split CP. & Can post-hoc sharpening match conformal retraining? \\
\textbf{B5. Post-hoc scores} & Apply APS, RAPS, and $-\log(T\pi_t)$ to B1 without retraining. & Can post-hoc score design close the gap? \\
\textbf{B6. Floors} & Uniform random scores; empirical (agent, position-bin) error prior in $\Dtr$. & Confirms encoder beats trivial priors. \\
\textbf{B7. External (opt.)} & Who\&When step-by-step LLM judge and OAT scores + vanilla split CP. & Context against published attributors. \\\bottomrule
\end{tabularx}}
\medskip

\begin{tcolorbox}[critbox,title={Claim Criteria (When WP1 Supports RQ1)}]
\small
(1) $\widehat{\mathrm{Cov}}$ lies inside the $\mathrm{Beta}(k,m+1-k)$ calibration interval on $\Dtest$ and recalibrated external splits.
(2) $\widehat{\mathrm{Ineff}}_{\mathrm{rel}}$ and $\widehat{\mathrm{Ineff}}$ are statistically lower than \textbf{B1--B5} at $\alpha=0.1$ without degrading at $\alpha\in\{0.05,0.2\}$.
(3) Gains are not driven by empty-set spikes or systematic under-coverage on long trajectories.
\end{tcolorbox}

% =====================================================================
\clearpage
\appendix
\section{Ablations}

\subsection{Alternative Conformal Training: ConfTr} \label{app:conftr}
As an alternative to DPSM's persistent threshold $q$ and fixed $D_1/D_2$ split, ConfTr (Stutz et al., ICLR 2022) simulates split CP inside every mini-batch via a differentiable quantile operator.

\begin{tcolorbox}[algbox,title={Algorithm 2 \textperiodcentered\ ConfTr Alternative (Differentiable In-Batch Calibration)}]
\small
\textbf{Inputs:} Full $\Dtr\setminus\Dval$; target $\alpha=0.1$; $\lambda, \tau, \varepsilon_{\mathrm{sort}}, \eta$; hinge $\kappa\in\{0,1\}$; batch size $b\in\{64,128\}$; steps $K$; $\varepsilon=10^{-6}$.
\begin{enumerate}[leftmargin=1.4em,itemsep=2pt]
\item \textbf{For} $k=1,\dots,K$:
  \begin{enumerate}[label=\arabic{enumi}.\arabic*.,leftmargin=1.6em,itemsep=1.5pt]
  \item Sample batch $B\subset\Dtr\setminus\Dval$ of size $b$ and split randomly into disjoint halves $B_{\mathrm{cal}}, B_{\mathrm{pred}}$ ($|B_{\mathrm{cal}}|=|B_{\mathrm{pred}}|=b/2$).
  \item \textbf{Differentiable calibration ($B_{\mathrm{cal}}$):} compute $\Smin_i(\theta)=1-\max_{t\in Y_i}\pi_t(x_i)$ for $i\in B_{\mathrm{cal}}$ and
  \[
  \tilde{q}_\theta = \mathrm{SoftQuantile}_{1-\alpha}\!\big(\{\Smin_i(\theta)\}_{i\in B_{\mathrm{cal}}};\,\varepsilon_{\mathrm{sort}}\big)
  \]
  via differentiable sorting (Blondel et al., 2020), keeping gradients through $\tilde{q}_\theta$.
  \item \textbf{Classification and relative soft set-size losses ($B_{\mathrm{pred}}$):} compute $\Lcls(\theta; B_{\mathrm{pred}})$ as in Algorithm~1 and
  \[
  \Lsize^{\mathrm{ConfTr}}(\theta) = \frac{1}{|B_{\mathrm{pred}}|}\sum_{j\in B_{\mathrm{pred}}}\frac{1}{T_j}\max\!\left(0,\; \sum_{t=1}^{T_j}\sigma\!\left(\frac{\tilde{q}_\theta - 1 + \pi_t(x_j)}{\tau}\right) - \kappa\right).
  \]
  \item \textbf{Encoder step:} $\theta \leftarrow \mathrm{AdamW}\big(\theta,\; \nabla_\theta[\Lcls(\theta; B_{\mathrm{pred}}) + \lambda\Lsize^{\mathrm{ConfTr}}(\theta)],\; \eta\big)$.
  \end{enumerate}
\item \textbf{Output:} Discard $\tilde{q}_\theta$; select best checkpoint $\theta^\star$ on $\Dval$ and calibrate $\Qhat(\alpha)$ on $\Dcal$ (Section~\ref{sec:vanilla-cp}).
\end{enumerate}
\end{tcolorbox}
\noindent Here $\kappa=0$ matches Algorithm~1 directly, while $\kappa=1$ is ConfTr's standard hinge discouraging empty sets. Compared to DPSM, ConfTr uses all of $\Dtr\setminus\Dval$ without a static split and eliminates $q$ and $\gamma$ (no pinball lag), at the cost of requiring larger batches ($|B_{\mathrm{cal}}|\ge 32$ at $\alpha=0.1$) and tuning $\varepsilon_{\mathrm{sort}}$.

\subsection{Score and Objective Alternatives}
\begin{itemize}
\item \textbf{Length-normalized log-prob score:} $S_\theta^{\mathrm{log}}(x,t) = -\log(T\pi_t(x)+\varepsilon)$ (or $-\log(\pi_t+\varepsilon)$). Removes $O(1/T)$ softmax scaling across variable $T$ and avoids $q\to 1$ compression when top-1 accuracy is $<50\%$ (analogous to ConfTr \texttt{THR LP}). Evaluated post-hoc in B5; directly swappable into Algorithm~1 (removing $[0,1]$ clipping on $q$) and Algorithm~2.
\item \textbf{Hard set-size objective:} $\frac{1}{|B^2|}\sum_{j}\frac{1}{T_j}\sum_{t}\ind[S_\theta(x_j,t)\le q]$. Eliminates $\tau$ and matches the exact metric, but has zero gradient w.r.t.\ $\pi_t$ almost everywhere, requiring straight-through or black-box gradient estimators for supervised training.
\end{itemize}

\subsection{Architecture and Backbone Ablations}
All ablations use activations from a single selected layer $\ell$ directly and keep the conformal training loop unchanged:
\begin{itemize}
\item \textbf{Step pooling:} Compare default mean pooling against \texttt{last\_token} pooling ($H_{\max I_t}^{(\ell)}$) and concatenated \texttt{[mean, last\_token]} pooling per step span $I_t$.
\item \textbf{Step-transition input:} Feed step deltas into the input projection, $u_t = \mathrm{LN}\big(W[h_t;\, h_t - h_{t-1}]\big)$ (with $h_0 = h_1$), to expose where the trace first diverges.
\item \textbf{Off-track auxiliary supervision:} Add a linear BCE head on $v_t$ predicting $z_t = \ind[t \ge \min Y]$ (whether the trajectory is already off-track at step $t$), trained on $B^1$ (or $B_{\mathrm{pred}}$) alongside $\Lcls$ with weight $\beta_{\mathrm{aux}}$ (no change to $B^2$ or $q$).
\item \textbf{Embedding backbones:} Swap frozen Qwen3-8B for \textbf{Qwen3-Embedding-\{0.6B, 4B, 8B\}} using the same Step~1 activation extraction to test dedicated embedding models and backbone scaling.
\end{itemize}

% =====================================================================
\clearpage
\section{Derivations and Caveats}

\subsection{Exact Gradients of the DPSM Objective}
Write $\ell_k=w^\top v_k$, $P_Y=\sum_{s\in Y}\pi_s$, $\tilde c_t=\sigma((q-1+\pi_t)/\tau)$, and $\tilde c'_t=\frac{1}{\tau}\tilde c_t(1-\tilde c_t)$. Using $\frac{\partial\pi_t}{\partial\ell_k}=\pi_t(\delta_{tk}-\pi_k)$:
\begin{align}
\frac{\partial \mathcal{L}^{\mathrm{cls}}}{\partial \ell_k} &= \frac{P_Y}{P_Y+\varepsilon}\,\pi_k - \frac{\ind[k\in Y]}{P_Y+\varepsilon}\,\pi_k, &
\frac{\partial \mathcal{L}^{\mathrm{size}}}{\partial \ell_k} &= \frac{\pi_k}{T}\Big(\tilde c'_k - \sum_{t=1}^T \pi_t\tilde c'_t\Big), &
\partial_q \rho_\alpha(q,S) &= \alpha - \ind[S\ge q]. \label{eq:wp1-grads}
\end{align}
$\tilde c'_k$ peaks near $\pi_k \approx 1-q$, focusing size gradients on borderline steps. The pinball step pushes $q$ up when batch miss rate exceeds $\alpha$ and down otherwise.

\subsection{Vanilla Split-CP Coverage Proof (Error-Set Labels)} \label{app:cp-proof}
For non-empty $Y\subseteq[T]$ and any $q\in\R\cup\{+\infty\}$, $\C_q(x)\cap Y\neq\emptyset \iff \exists t\in Y: S_\theta(x,t)\le q \iff \Smin_\theta(x,Y)\le q$. Conditioned on $\theta$ (independent of $\Dcal$ and test $(x,Y)$), applying $(x,Y)\mapsto\Smin_\theta(x,Y)$ preserves exchangeability. By the split-CP rank lemma with $k=\lceil(1-\alpha)(m+1)\rceil$,
\[
1-\alpha \le \frac{k}{m+1} \le \Prob\big(\C_\theta(x)\cap Y\neq\emptyset\big) < 1-\alpha + \frac{1}{m+1},
\]
(upper bound requires distinct scores). Conditional on $\Dcal$, test coverage follows $\mathrm{Beta}(k,m+1-k)$ (std $\approx\sqrt{\alpha(1-\alpha)/m}$). The proof makes no assumption on $T$, step dependence, or scorer accuracy.

\subsection{Caveats and Diagnostics}
\begin{itemize}
\item \textbf{Empty-set gaming:} HPS allows $\C(x)=\emptyset$, so marginal size minimization can drop the hardest/longest $\alpha$ traces. Check empty-set rate and length-tercile coverage.
\item \textbf{Pinball lag:} Starting at $q_0=0.5$ without warm-up causes early lag; log $q$ vs.\ exact $D_1$ quantile each epoch to verify $\gamma$ tracks $\theta$.
\item \textbf{Sample-size split:} Compare DPSM against both B1 (full $\Dtr$) and B2 ($D_1$ only) to separate loss benefits from the $50\%$ label split.
\item \textbf{Frozen-feature ceiling:} ConfTr reported modest THR gains ($2$--$3\%$) on frozen features; comparing WP1 vs.\ WP2 isolates the gain of adapting representations.
\item \textbf{Transfer shift:} Benchmarks differ in label semantics ($|Y|=1$ vs.\ multi-step sets); zero-shot transfer coverage reflects $P_{Y|X}$ shift, so use recalibrated splits for efficiency comparisons.
\end{itemize}

\end{document}
